\section{Generalization with Terminal State Labels}
\label{sec:lengthgen-arch}

We separate supervised extrapolation from terminal-reward optimization. Each
model observes a binary event sequence and receives a single label for the final
residue. Every model minimizes terminal-label NLL for 22,000 Adam updates at
learning rate $0.0015$, batch size 256, and lengths $4,6,\ldots,14$. The archived
phase-rotation configuration uses an angle lattice and Fourier readout initialization.
The comparison varies architectural priors, including the structure supplied
to the phase-rotation model.

Table~\ref{tab:lengthgen-arch} compares ten seeds at the common horizon $L=400$.
The phase rotation with an angle lattice solves a majority of seeds at every
cycle size. GRU and LSTM nearly always extrapolate at $k=3$--$5$ in this
terminal-label protocol. The successful-seed counts at $L=1600$ for
$k=2,3,4,5$ are
$(9,9,6,9)$ for the phase-rotation model, $(0,10,10,9)$ for GRU, and $(1,9,10,10)$ for LSTM.
The archived predictions reproduce the counts.
\begin{table}[H]
\centering\small
\caption{Extrapolation after terminal-label supervision. Entries count
runs exceeding $0.90$ greedy accuracy at $L=400$ out of ten seeds per
configuration. Recurrent models use $3\times512$ evaluation sequences per
seed, and the Transformer uses 128.}
\label{tab:lengthgen-arch}
\begin{tabular}{lccccc}
\toprule
Architecture & Parameters & $\Z_2$ & $\Z_3$ & $\Z_4$ & $\Z_5$ \\
\midrule
Phase rotation (angle lattice, Fourier initialization) & 337--433 & 9/10 & 9/10 & 6/10 & 9/10 \\
GRU & 3.7--3.8k & 0/10 & 10/10 & 10/10 & 10/10 \\
LSTM & 5.0--5.1k & 1/10 & 10/10 & 10/10 & 10/10 \\
Transformer (no positional encoding) & 20k & 0/10 & 0/10 & 0/10 & 0/10 \\
\bottomrule
\end{tabular}
\end{table}

